\section{Methode 1: Risikobeurteilung}
Die Risikobeurteilung (auch Risk-Assessment genannt) ist ein zentraler Bereich des Risikomanagement Werkzeugkastens. Ziel ist es Risiken zu identifizieren und hinsichtlich ihrer
qualitativen und quantitativen Folgen zu analysieren \parencite{walder_risikomanagement_2017}. 
Das Ergebnis daraus ist die Bewertung, die für alle weiteren Abwägungen genutzt werden kann.

Dafür sollen in diesem Kapitel unterschiedliche Methoden der Risikoidentifikation und -bewertung aufgeführt werden, die später Anwendung in 
der Risikomanagement Strategie in Data-Science Projekten finden kann.

\subsection{Identifikation}
Zur Identifikation von Risiken bieten sich einige verschiedene Methoden an. \parencite{blender_risiken_2025}
\begin{itemize}
    \item Brainstorming - Kollaborativer Austausch im Team.
    \item SWOT-Analyse - Strukturierte Stärken und Schwächen Analyse.
    \item Experten-Interviews - Experte aus der Fachdomäne zu Risiken befragen.
    \item Vergangenheitsdaten - In der Vergangenheit aufgetretene Schäden hinsichtlich ihrer Ursachen analysieren.
    \item Mitarbeiter-Interviews - Befragungen in der eigenen Organsiation durchführen und dort Meinungen und Bedenken einholen.
\end{itemize}
Diese Methoden beziehen sich im allgemeinen ehr auf eine klassische Projektorganisation und werden den gesonderten Bereichen der 
Data-Science nicht vollumfänglich gerecht. Wenn zum Beispiel die Bereiche der \textbf{Datenbezogenen-} oder \textbf{Modell-Risiken} betrachtet werden 
könnten die bisher aufgeführten Identifikationsmethoden nicht die nötige Tiefe liefern, die es an dieser Stelle braucht.
Auch der Bereich der \textbf{Strategischen Risiken} ist bei Data-Science Projekten geprägt von Besonderheiten. 
IT-Infrastruktur, eine breit zu behandelnde Fachdomäne und technische Hürden erfordern ein geschäftes Augenmaß auf 
der strategischen Ebene um Risiken adäquat zu identifizieren und zu bewerten.

(Herleiten der Data-Science spezifischen Methoden) \dots

Für die Datenbezogenen und Modell Risiken gibt es aus der Praxis bereits einige erprobte Methoden die z.B. Probleme 
wie Model-Drift, Data-Drift, Ethic-Biases bereits frühzeitig identifiziert und Zeit lässt zu agieren.

(Methoden auflisten) \dots

\subsection{Analyse und Bewertung}

(Einige der Datenspezifischen Methoden die bereits bei der Identifikation genannt wurden, beziehen die Analyse und Herkunft des Problems bereits mit ein.) \dots

(Für die Bewertung wird man hier sicherlich wieder auf die gängigen Riskmanagement Methoden zurückgreifen. Sprich - Risikomatrix, Assessment Report, Risk Register, usw.)  \dots

Qualitative Methoden befassen sich mit der Auflistung und Einordnung ohne Auswertung. z.B. Lieferengpässe (Wahrscheinlichkeit: Hoch, Auswirkung: Mittel).
Quantitative Methoden werten diese Auflistung aus und bieten eine art Score, der bei der Abwägung hilft. 
z.B. Lieferengpässe (Wahrscheinlichkeit: 7/10, Auswirkung: 6/10 - Riskscore: 7 x 6 = 42)

(Herleiten und mit Quellen belegen)
\subsubsection{Qualitative Methoden}
\begin{itemize}
    \item Risikomatrix
    \item Risikokategorisierung
    \item MoSCoW Methode
\end{itemize}
Die Methoden unterscheiden sich lediglich hinsichtlich ihrer Gruppierung der Risiken.

\subsubsection{Semi-quantitative Methoden}
\begin{itemize}
    \item Risikoscore
    \item FMEA
\end{itemize}

\subsubsection{Quantitative Methoden}

\subsection{Beispiel-Artefakt: Risk Register}
\begin{longtable}{@{}p{3.0cm}p{3.2cm}p{2.1cm}p{2.1cm}p{4.2cm}@{}}
\toprule
\textbf{Risiko} & \textbf{Ursache/\allowbreak Trigger} & \textbf{Eintritt} & \textbf{Aus\allowbreak wirkung} & \textbf{Maßnahmen /\allowbreak Owner} \\
\midrule
\endfirsthead
\toprule
\textbf{Risiko} & \textbf{Ursache/\allowbreak Trigger} & \textbf{Eintritt} & \textbf{Aus\allowbreak wirkung} & \textbf{Maßnahmen /\allowbreak Owner} \\
\midrule
\endhead
\textless Risiko 1\textgreater & \textless Trigger\textgreater & H/M/L & H/M/L & \textless Maßnahme + Verantwortliche:r\textgreater \\
\textless Risiko 2\textgreater & \textless Trigger\textgreater & H/M/L & H/M/L & \textless Maßnahme + Verantwortliche:r\textgreater \\
\bottomrule
\end{longtable}
